\chapter{Design}\label{ch:design}

This chapter presents DUET's architecture, dual-index warning scheme, windowed consistency proof, and
instantiation of ACKA. \Cref{app:algorithms} gives the pseudocode, and \cref{ch:eval} evaluates the
associated costs.

\paragraph{Notation.} Hashing is domain-separated in the manner of RFC~6962
(\cref{ssec:bg-merkle}): $\Hash_{\mathsf{leaf}}(x) = \Hash(\mathtt{0x00} \concat x)$
and $\Hash_{\mathsf{node}}(l, r) = \Hash(\mathtt{0x01} \concat l \concat r)$, with
$\Hash$ instantiated as SHA-256 and $\concat$ denoting byte concatenation. The empty
position has the constant default hash $\Hash(\varepsilon)$. The prefix tree and the log
use these same leaf and node tags. Cross-tree confusion is not exploitable here, since
prefix roots appear only as log leaf data and log roots only in signed tree heads, but a
production deployment should apply distinct per-structure tags as defence in depth.
Identities are treated as fixed-width service identifiers, so concatenations of identity and
label material are unambiguous without explicit length prefixes.

% ============================================================================
\section{Architecture}\label{sec:design-arch}

The proposed architecture is built around three authenticated data structures arranged in layers,
where each layer authenticates the one below (\cref{fig:design-arch}). The prefix tree stores the
current directory state, the append-only log records every committed version of that state, and the
signed tree head authenticates the latest log root. This is the combined-tree design of the IETF Key
Transparency working group (\cref{ssec:bg-keytrans}). Lookups are served against the \emph{last
committed epoch}, never the live tree. Only a committed prefix root has been appended to the log and
can therefore be linked by an inclusion proof to a signed tree head.

\begin{figure}[t]
\centering
% fig-design-arch.tex — extracted from ch4-design.tex
% Label: \label{fig:design-arch}
\begin{tikzpicture}[>=Stealth,font=\small,
   layer/.style={draw,rounded corners,minimum width=6.4cm,minimum height=11mm,align=center}]
  \node[layer] (sth) at (0,4) {Signed Tree Head (STH)\\[-1pt]\footnotesize signed commitment};
  \node[layer] (log) at (0,2) {Append-only Merkle log\\[-1pt]\footnotesize history of directory states};
  \node[layer] (pt)  at (0,0) {Sparse Merkle prefix tree\\[-1pt]\footnotesize current authenticated directory};
  \draw[->] (pt) -- (log);
  \draw[->] (log) -- (sth);
  \node[draw,rounded corners,align=center,font=\footnotesize] (vrf) at (5.2,0) {VRF\\global labels};
  \draw[->,dashed] (vrf) -- (pt);
\end{tikzpicture}

\caption[The authenticated directory core.]{The authenticated directory core. The prefix tree,
Merkle log, and signed tree head form the commitment chain. The VRF derives global directory
labels.}
\label{fig:design-arch}
\end{figure}

\subsection{The Three Structures}\label{ssec:design-structures}

The three structures form a verification stack. A lookup proof is checked against the prefix root,
the prefix root's inclusion in the Merkle log, and the signed tree head that authenticates the log
root (\cref{ssec:design-verify}).

\paragraph{Prefix Tree.} The directory is a sparse binary Merkle tree that maps each
256-bit \emph{label} to a \emph{value}. For a user-identity entry, the value is an
identity public key. The label is the position of the entry in the tree, derived
rather than chosen, and it comes in two families: a \emph{global slot}, computed
server-side as a VRF over the user identifier, and the client-derived \emph{raw slots},
computed from conversation secrets: the per-conversation dual-index lookup slots and shared
warning slot (\cref{sec:design-dual}) and, in continuous mode, the per-epoch CEA slots
(\cref{ssec:design-modes}).

\paragraph{Append-Only Log.} Each committed prefix-tree root is appended to a Merkle
log (\cref{ssec:bg-merkle}) as a single leaf $\prefroot{t}$. The log therefore records
every committed directory state, and a consistency proof certifies that this
history has only ever grown. Nothing previously committed has been modified,
deleted, or reordered.

\paragraph{Signed Tree Head (STH).} Each epoch, the root of the log is signed with Ed25519
and published as the STH. This gives clients and auditors a single authenticated value
that commits to the whole directory and its history. It is the anchor against which
every lookup is ultimately checked.

\paragraph{Why a VRF?} The VRF carries G3 (Directory Privacy)
(\cref{sec:threat-goals}). The operator holds the VRF secret key and maps a user
identifier to a pseudorandom label, attaching a proof that the label was computed
correctly. A client verifies that proof against the operator's public key but cannot
itself compute or predict any other label (\cref{ssec:bg-ct2kt}). The labels are thus
deterministic for the operator yet unpredictable to outsiders, which prevents an
observer from enumerating the directory.

\subsection{Epoch Lifecycle: Update, Commit, and the Signed Tree Head}\label{ssec:design-epoch}

The directory evolves in epochs. An $\textsf{update}(\mathit{label}, \mathit{value})$
changes the value stored at one leaf of the current uncommitted prefix tree and recomputes only the
hashes on the path from that leaf to the root. The running time is therefore $O(256)$, one hash per level of the fixed 256-bit label space
and so constant in $N$, rather than $O(N)$. The number of \emph{transmitted} non-default siblings in the resulting proof is a
separate quantity, whose expected value is $\approx\lceil\lg N\rceil$.

A $\textsf{commit\_epoch}()$ then closes the epoch. It reads the current prefix-tree
root, appends it to the log as the next leaf, recomputes the log root, snapshots the
committed tree so that subsequent searches are served against a frozen state, and signs
the new head.

The contents of the signed tree head determine exactly what the signature
authenticates. The STH signs, under Ed25519, the 88-byte little-endian message
\begin{equation}\label{eq:design-sth}
  \mathit{epoch} \concat \mathit{tree\_size} \concat \mathit{root} \concat
  \mathit{timestamp} \concat \Hash(\mathit{config}),
\end{equation}
where $\Hash(\mathit{config}) = \Hash(\mathit{pk}_{\mathsf{vrf}} \concat
\mathit{pk}_{\mathsf{sth}} \concat \mathit{suite} \concat \mathit{version} \concat
\mathit{epoch\_interval} \concat \mathit{raw\_label\_policy})$, the trailing byte binding the
write-once raw-label policy so the server cannot silently relax it. Each field of \cref{eq:design-sth} is necessary for a different aspect of the verification procedure. The $\mathit{tree\_size}$ is
required by the inclusion and consistency verifiers. Signing the $\mathit{timestamp}$
means the timestamp cannot be altered without invalidating the signature. Enforcing \emph{freshness}
additionally requires a client-side maximum-age or publication-cadence policy, which the prototype
does not implement (\cref{ch:discussion}). The trailing $\Hash(\mathit{config})$ is a
CONIKS-style policy commitment (the field CONIKS calls $P$) that binds the provider's
parameters into the signature. With it, the server cannot silently swap its VRF key,
its signing key, the ciphersuite, the protocol version, or the epoch cadence, or
serve different parameters to different clients, because any client that recomputes
$\Hash(\mathit{config})$ from the parameters it expects will reject the mismatch. A
parameter substitution thus becomes a signed, monitorable act rather than a silent one: a uniform
mismatch is rejected by any client that recomputes the commitment, and serving different parameter
sets to different clients is equivocation on the configuration commitment, exposed by the same
cross-vantage comparison as A4.

\subsection{The Client Verification Chain}\label{ssec:design-verify}

A client accepts no server response without independently verifying it: every lookup is
ultimately authenticated against a signed tree head. Each verification traverses the same three-layer architecture
introduced in \cref{sec:design-arch}: first the prefix tree, then the append-only log,
and finally the signed tree head. The verification chain is the six-step pipeline of
\cref{fig:design-verify}.

\begin{figure}[t]
\centering
% fig-design-verify.tex — the client verification chain (VerifiedLookup).
% Label: \label{fig:design-verify}
\begin{tikzpicture}[>=Stealth,font=\footnotesize,
   vstep/.style={draw,rounded corners,text width=13cm,minimum height=8mm,align=left,inner sep=5pt},
   note/.style={font=\scriptsize\itshape,text=black!55,text width=13cm,align=left}]
  \node[vstep] (s1) {\textbf{1. Requested-label binding} (every path): require $\pi_{\mathrm{prefix}}.\mathit{label}$ equals the label the client queried, so a valid proof for a different label is rejected};
  \node[vstep,below=3mm of s1.south,anchor=north] (s2) {\textbf{2. VRF label binding} (VRF path only): verify the ECVRF proof and require $\beta[0{:}32]$ equals that same queried label (RFC~9381)};
  \node[vstep,below=3mm of s2.south,anchor=north] (s3) {\textbf{3. STH signature and log history}: verify the Ed25519 signature and require the returned head to match the caller-supplied verified head. Between lookups, the Signal Desktop verifier accepts a newer head only after a valid consistency check from the last trusted head};
  \node[vstep,below=3mm of s3.south,anchor=north] (s4) {\textbf{4. Value binding}: require the prefix proof binds to the returned value};
  \node[vstep,below=3mm of s4.south,anchor=north] (s5) {\textbf{5. Prefix root}: recompute the prefix-tree root from the proof (siblings $+$ 32-byte bitmap)};
  \node[vstep,below=3mm of s5.south,anchor=north] (s6) {\textbf{6. Log inclusion and latest-state binding}: check the recomputed root is a leaf under the signed log root, and that the inclusion covers the \emph{last} log leaf (inclusion of an older leaf does not establish the current directory state)};
  \foreach \a/\b in {s1/s2,s2/s3,s3/s4,s4/s5,s5/s6}{\draw[->] (\a) -- (\b);}
  \node[note,below=3mm of s6.south,anchor=north] {The warning slot is checked by a separate invocation of this same chain (\textsf{CheckWarning}), so its verified inclusion or non-inclusion is established independently.};
\end{tikzpicture}

\caption[The six-step client verification chain.]{The six-step client verification chain. Every lookup is checked locally
against a signed root.}
\label{fig:design-verify}
\end{figure}

The verification chain has three properties that the figure does not make evident. First, global
lookups require the step~2 VRF check, while client-derived raw labels omit it, and there the step~1
requested-label binding stands on its own. Second, the signed head used for verification must be
identical to the head the client independently obtained and verified, either the latest head it
fetched or, during catch-up, the new head whose append-only consistency with its last trusted head it
has just checked. The client rejects any response anchored to a different head, even one that carries
a valid signature, which prevents stale-snapshot acceptance and prevents an equivocating server from
substituting another signed view. Third, inclusion is pinned to the latest leaf, because each commit
appends the new directory root as the log's last leaf, so step~6 requires
$\mathit{leaf\_index} = \mathit{tree\_size} - 1$ at exactly the signed tree size, and a server
therefore cannot answer with a stale directory snapshot under the current signed head. Inclusion
proves membership, not \emph{currency} (RFC~9162~\cite[Section~2.1.3.2]{rfc9162}), and the currency
pin closes that gap.

The Rust and TypeScript implementations enforce the same six-check chain, and \cref{ch:eval} measures
the full chain. The ``full client verification'' benchmarked in \cref{ch:eval} is the complete
six-step chain.

% ============================================================================
\section{The Dual-Index Warning Scheme}\label{sec:design-dual}

The dual-index scheme meets the in-band peer-warning requirement of \cref{ssec:bg-implications} and so
achieves G1 (Key-Substitution Detection) against the substitution attacks A1, A3, and A5. An in-band
warning channel flags a server-side slot substitution (A1) and its continuous-mode counterpart (A5).
The secret-derived dual-index labels detect a network man-in-the-middle (A3), because an attacker with
a different shared secret derives disjoint labels, so the proof-verified lookup returns non-inclusion.
A comparison with the peer's global VRF-anchored directory entry also detects pre-emptive occupation
of the expected slot, because an attacker-controlled value cannot match the authenticated global key.

\subsection{Label Derivation}\label{ssec:design-labels}

The dual-index scheme allocates three directory slots to each conversation: two lookup
slots (one per party) and one shared warning slot (\cref{fig:design-dualindex}).
Alice and Bob derive these slot labels from the X3DH/PQXDH shared secret $\sharedS$
(\cref{ssec:bg-x3dh}) using HKDF-SHA256 as specified in RFC~5869~\cite{rfc5869}, the same
key-derivation function Signal already uses elsewhere:
\begin{equation}\label{eq:design-labels}
\begin{aligned}
  \mathit{PRK} &= \HKDFex(\text{salt} = \text{\texttt{"duet-dual-index"}},\ \sharedS)\\
  \ell_{\mathsf{self}} &= \HKDFexp\!\big(\mathit{PRK},\ \text{\texttt{"lookup"}} \concat \mathit{id}_{\mathsf{self}} \concat \mathit{id}_{\mathsf{peer}},\ 32\big)\\
  \ell_{\mathsf{peer}} &= \HKDFexp\!\big(\mathit{PRK},\ \text{\texttt{"lookup"}} \concat \mathit{id}_{\mathsf{peer}} \concat \mathit{id}_{\mathsf{self}},\ 32\big)\\
  \ell_{\mathsf{warn}} &= \HKDFexp\!\big(\mathit{PRK},\ \text{\texttt{"warning"}} \concat \min(\mathit{id}_\A, \mathit{id}_\B) \concat \max(\mathit{id}_\A, \mathit{id}_\B),\ 32\big).
\end{aligned}
\end{equation}
In \cref{eq:design-labels}, $\mathit{PRK}$ is the extracted pseudorandom key, while
$\mathit{id}_{\mathsf{self}}$ and $\mathit{id}_{\mathsf{peer}}$ identify the deriving party and its
peer. The $\min$ and $\max$ operations order the fixed-width identity encodings lexicographically.
Each output is 32 bytes, matching the prefix tree's label width. Canonical ordering is used only for
the warning label, which both parties must derive identically. Directional ordering gives each party a
separate lookup slot.

\begin{figure}[ht]
\centering
% fig-design-dualindex.tex — extracted from ch4-design.tex
% Label: \label{fig:design-dualindex}
\begin{tikzpicture}[font=\small,>=Stealth,
   slot/.style={draw,rounded corners,minimum width=2.6cm,minimum height=10mm,align=center},
   sec/.style={draw,rounded corners,fill=black!6,minimum width=1.9cm,minimum height=13mm,align=center},
   kdf/.style={draw,rounded corners,fill=black!3,align=center,minimum height=12mm},
   elbl/.style={font=\scriptsize,align=center}]
  \node[sec] (S) at (0,0) {$\sharedS$\\[-1pt]\footnotesize shared\\[-3pt]\footnotesize secret};
  \node[kdf,minimum width=2.9cm] (ext) at (3.1,0)
    {HKDF-Extract\\[1pt]{\scriptsize salt ``\texttt{duet-dual-index}''}};
  \node[draw,rounded corners,minimum width=1.0cm,minimum height=8mm] (prk) at (5.7,0) {PRK};
  \node[slot] (self) at (10.7,1.95)  {$\ell_{\mathsf{self}}$\\[-1pt]\footnotesize Alice's key};
  \node[slot] (peer) at (10.7,0)     {$\ell_{\mathsf{peer}}$\\[-1pt]\footnotesize Bob's key};
  \node[slot,fill=black!10] (warn) at (10.7,-1.95) {$\ell_{\mathsf{warn}}$\\[-1pt]\footnotesize shared tripwire};
  \draw[->] (S.east) -- (ext.west);
  \draw[->] (ext.east) -- (prk.west);
  \draw[->] (prk.east) -- node[elbl,above,pos=0.55]{\texttt{Expand}\\``lookup''$\,\|\,\mathit{id_A}\|\mathit{id_B}$} (self.west);
  \draw[->] (prk.east) -- node[elbl,above]{\texttt{Expand}\\``lookup''$\,\|\,\mathit{id_B}\|\mathit{id_A}$} (peer.west);
  \draw[->] (prk.east) -- node[elbl,below,pos=0.55]{\texttt{Expand}\\``warning''$\,\|\,\min\|\max$} (warn.west);
\end{tikzpicture}

\caption[The dual-index scheme's three derived slots.]{Two lookup slots and a shared warning slot
are derived from the conversation secret $\sharedS$. HKDF-Extract produces $\mathit{PRK}$.
Three domain-separated HKDF-Expand operations then produce $\ell_{\mathsf{self}}$,
$\ell_{\mathsf{peer}}$, and $\ell_{\mathsf{warn}}$. Alice and Bob share $\sharedS$, so they derive
the same warning label and can read either party's warning. A different secret $\sharedS'$ produces
a disjoint label set.}
\label{fig:design-dualindex}
\end{figure}
\FloatBarrier

This corrects two simplifications that appeared in early descriptions of the scheme: it
uses HKDF rather than a bare hash of $\sharedS$, and it uses three labels with canonical
ordering applied to the warning label alone. Cross-implementation test vectors verify the
byte-for-byte equivalence of the derivation across the Rust and TypeScript implementations on the
committed cross-language vectors (\cref{sec:tests-xlang}).

\subsection{The Three Security Properties}\label{ssec:design-properties}

The derivation establishes three properties on which detection rests.

\paragraph{Symmetry.} From Alice's and Bob's respective viewpoints, the directional lookup
derivation ensures $A.\ell_{\mathsf{self}} = B.\ell_{\mathsf{peer}}$ and $A.\ell_{\mathsf{peer}} =
B.\ell_{\mathsf{self}}$, while the canonical identity ordering gives $A.\ell_{\mathsf{warn}} =
B.\ell_{\mathsf{warn}}$. The two parties therefore agree on which slot each key lives in and, above
all, on the one warning slot they both watch.

\paragraph{Pre-Use Unpredictability.} The server does not hold $\sharedS$, so the HKDF outputs are
pseudorandom from the server's perspective and it cannot distinguish $\ell_{\mathsf{warn}}$ from
any other leaf in the tree. It can therefore neither pre-position a value on the warning slot nor
recognize and selectively censor it. This protection holds until the label is first queried. After
first use the operator can associate a conversation's labels, and those consequences are handled by
the write-once policy, the transition audit, and the persistence check (\cref{ssec:design-detection}).

\paragraph{Scope of the Warning Channel.} The shared warning channel applies only after an honest
X3DH/PQXDH handshake has established a common secret. It covers A1 and A5 for established pairs.
First-contact A3 is detected through divergent labels and comparison with the peer's global VRF-bound
entry, while a substitution that predates an honest baseline remains the TOFU boundary
(\cref{sec:threat-fair}).

\subsection{Detection and the Bidirectional Warning}\label{ssec:design-detection}

On detecting a suspicious key at a lookup slot, a party writes data to
$\ell_{\mathsf{warn}}$. Both parties' background monitors read this slot on their next check and
raise a warning. The payload carries a MAC under a key derived from the shared secret and the
canonically ordered identities
($\HKDFexp(\mathit{PRK}, \text{\texttt{"warn-mac"}} \concat \min(\mathit{id_A},\mathit{id_B}) \concat \max(\mathit{id_A},\mathit{id_B}), 32)$).
Only the two conversation parties can produce or verify this MAC. The warning label is unpredictable
to the operator until one party reveals it in a request, so any \emph{unauthenticated} value returned
at that location is itself suspicious and still raises a warning. After the operator observes the
label, the MAC allows the peer to distinguish an authentic warning from an operator-injected value
(\cref{ssec:design-properties}).

An append-only log authenticates committed roots but does not by itself constrain directory
transitions. DUET therefore combines write-once labels, per-epoch transition auditing, and retained
persistence evidence. \Cref{sec:sec-suppression} gives the full suppression and offline-detection
analysis, including the A2 compound condition and the treatment of refused responses.

\subsection{Shared-Label Consistency}\label{ssec:design-consistency}

The dual-index scheme's detection rests on a static, same-head consistency statement.
Under one signed tree head the label Alice reads for Bob's key and the label Bob self-checks are the
same label, so it cannot carry two different authenticated values. Separate mechanisms address
warning publication, multiple epochs, offline catch-up, and equivocation. Because both parties derive
the same directional lookup label, the lemma holds even when $\sharedS$ is known.

\Cref{lem:dualindex-consistency} states this formally with its argument. Under one signed head a
substituted peer value and an honest self-check at the shared label cannot both verify, so a malicious
server cannot keep both consistent. Its interpretation and the surrounding detection cases are
developed in \cref{def:g1-consistency}. A complete formal security proof remains future work
(\cref{ch:discussion}).

The dual-index mechanism complements rather than replaces CEA. $\CEA$ authenticates the
values stored at directory entries by deriving per-epoch fingerprints. The dual-index
scheme instead derives the directory labels from the shared secret, identifying the
directory locations at which those values are stored. The two mechanisms therefore address complementary aspects of
authentication: $\CEA$ authenticates entry contents, while the dual-index scheme
authenticates entry placement.

\subsection{Session and Continuous Authentication Modes}\label{ssec:design-modes}

The scheme described so far is the \emph{session mode}: the labels come from the
\emph{initial} X3DH secret $\sharedS$, the registered value is the \emph{identity key},
and a party re-registers only when its signed pre-key rotates. Session mode detects the
handshake-time attacks (server substitution (A1) and network man-in-the-middle (A3))
at the cost of one log write per pre-key rotation.

The directory logs the long-term identity key. The remaining X3DH secrets contribute to the shared
secret $\sharedS$ from which the labels derive after honest establishment, rather than being logged as
individual objects. \Cref{tab:design-x3dh-coverage} sets this coverage against Signal's manual Safety
Number.

\begin{table}[H]
\centering
\caption[X3DH key-material coverage: Signal versus this thesis.]{Coverage of X3DH key material by
Signal's Safety Number and the scheme proposed in this thesis. The Signal column describes the Safety
Number. Optional AKV is discussed in \cref{ssec:bg-gap}.}
\label{tab:design-x3dh-coverage}
\small
\renewcommand{\arraystretch}{1.4}
\vspace{4pt}
\begin{tabularx}{\textwidth}{@{}>{\raggedright\arraybackslash}p{3.4cm}>{\raggedright\arraybackslash}X>{\raggedright\arraybackslash}X@{}}
\toprule
Cryptographic material & Signal's Safety Number & Proposed scheme (this work) \\
\midrule
Long-term identity key ($\pkid{}$) & Yes, by manual out-of-band comparison & Yes:
  automatically logged and checked by a background monitor \\
Signed pre-key ($\spk{}$) & No & Not logged individually. A signed-pre-key
  rotation triggers re-registration of the identity-key binding. The signed pre-key itself is not a
  logged object, its authenticity resting on its X3DH signature under the verified identity key \\
One-time pre-key ($\opk{}{i}$) & No & Not logged individually (single-use),
  contributes to $\sharedS$ \\
Ephemeral key ($\ek{}$) & No & Not independently logged, contributes to $\sharedS$, from which the
  conversation's labels derive after honest establishment \\
\bottomrule
\end{tabularx}
\end{table}
\FloatBarrier

Alongside it, the system implements a \emph{continuous mode} that binds authentication to
the \emph{entire} continuous key agreement, following ACKA's example CEA construction
(\cref{sec:bg-acka}). Each ratchet epoch carries a \emph{public} ratchet key $T_t$ and a \emph{secret} epoch key
$I_t$ (ACKA's Signal mapping sets $I_t = ck$, Signal's chain key, itself a KDF
of the ratchet Diffie--Hellman output, whereas this PoC mixes in the raw DH output as a stand-in for
$ck$, see \cref{ch:discussion}). At epoch $t$ the party
computes, from the current chain state $k_t$ (with $k_0$ the state established by Setup), an evolving
label and a fingerprint that authenticates $T_t$, and only then folds the secret $I_t$ into the
chain to produce $k_{t+1}$:
\begin{equation}\label{eq:design-cea}
\begin{aligned}
  L_t &= \Hash\!\big(k_t \concat \min(\mathit{id}_\A, \mathit{id}_\B) \concat \max(\mathit{id}_\A, \mathit{id}_\B) \concat t\big),\\
  F_t &= \Hash\!\big(\KDF(k_t,\ \mathit{id}_{\mathsf{sender}} \concat \mathit{id}_{\mathsf{peer}} \concat t) \concat T_t\big),\\
  k_{t+1} &= \KDF\!\big(k_t,\ \KDF(I_t, \text{\texttt{"CEA"}})\big) \qquad (\KDF = \text{HMAC-SHA256, a PRF}).
\end{aligned}
\end{equation}
It registers the pair $(L_t \to F_t)$ of \cref{eq:design-cea} in the same log, one entry per ratchet epoch. The
label $L_t$ uses the canonical $\min/\max$ identity ordering (the same convention as the warning
label in \cref{eq:design-labels}), so both parties derive the identical label for an epoch and it
maps to a single log slot. The fingerprint $F_t$, by contrast, is \emph{directional}
($\mathit{id}_{\mathsf{sender}} \concat \mathit{id}_{\mathsf{peer}}$): the verifier recomputes it with the
party who sent that epoch's ratchet key first, which binds the fingerprint to the sender's public
key $T_t$ and lets a substituted key be detected at the shared label.

The ordering of operations within an epoch follows ACKA's composed construction
(their Figure~8~\cite{dowling2023acka}) in three respects, each of which matters for correctness. First, the label and
fingerprint are computed over the \emph{current} (pre-update) state and the \emph{public} key, and
only then is the \emph{secret} $I_t$ mixed into the chain. Mixing the secret epoch key rather than
the public ratchet key keeps the chain unpredictable to an adversary that has compromised the
pre-shared state, the property the pre-use unpredictability and forgery-resistance arguments need. Before the
secret is folded into the CEA chain, it is domain-separated as $\KDF(I_t, \text{\texttt{"CEA"}})$
(ACKA Figure~8). In a faithful Signal integration $I_t$ is Signal's chain key $ck$, which also feeds the
session-key derivation (the paper's $\KDF(I_t, \text{\texttt{"CKA"}})$). The distinct \texttt{"CEA"}
and \texttt{"CKA"} tags keep the transparency-chain value independent of the session key, the key
separation the security argument relies on. Second,
before posting, the sender checks that the label $L_t$ is unoccupied: ACKA specifies that ``the
sender will also check if the label has already been uploaded to the log with a fingerprint, which
would indicate an impersonation''~\cite{dowling2023acka}, and the sender advances its chain only
after its write is confirmed in a committed head. Third, the receiver advances its chain only when a
verification returns \textsc{Consistent}. A \textsc{Fork} verdict leaves the chain unadvanced over the
disputed epoch and reports an in-band warning. The prototype does not block session establishment in
Signal. A pending lookup (the sender's fingerprint not yet committed) leaves the receiver's state
untouched, so the check can be retried without desynchronizing the chain.

The two designs represent label history differently. ACKA's log stores a vector of fingerprints for
each label. Its receiver accepts only a single fingerprint and treats two fingerprints at the same
label as evidence of impersonation. DUET stores one value at each label, so it implements the same
detection objective through two checks. If the adversary posts first, the honest sender's pre-post
check finds the label occupied and rejects the write. If the adversary later attempts to replace a
committed fingerprint, the transition audit rejects the update. Since each leaf hash records whether
the leaf is mutable or immutable, the auditor can identify and reject a rotation of an immutable
label. The receiver's per-epoch check cannot detect this overwrite when the adversary controls the
pre-shared state, so this case depends on the transition audit. Proving that ACKA's security argument
carries over to this directory representation remains future work (\cref{ch:discussion}).

For per-epoch CEA verification, detection follows ACKA. An active man-in-the-middle that injects a
forged ratchet key $T'$ either posts under a label the impersonated party never derives, since a
different chain produces a different $L_t$, or causes the committed fingerprint to differ from the
verifier's expected fingerprint at the shared label. In the latter case, \textsf{CeaVerify} returns
\textsc{Fork}, leaves the chain unchanged, and writes to the shared warning slot. This CEA verdict is
distinct from split-view equivocation (A4), which requires the cross-vantage auditor.
\Cref{fig:design-modes} contrasts session and continuous modes.

\begin{figure}[H]
\centering
% fig-design-modes.tex — extracted from ch4-design.tex
% Label: \label{fig:design-modes}
\begin{tikzpicture}[>=Stealth,font=\footnotesize,
   box/.style={draw,rounded corners,align=center,minimum width=2cm,minimum height=7mm}]
  \node[font=\small] at (0,1.2) {\textbf{Session mode}};
  \node[box] (s)  at (0,0.4)  {$\sharedS$};
  \node[box] (sl) at (0,-0.7) {labels};
  \node[box] (sk) at (0,-1.8) {identity key};
  \draw[->] (s) -- node[right,font=\footnotesize]{HKDF} (sl);
  \draw[->,dashed] (sl) -- node[right,font=\footnotesize]{stores} (sk);
  \node[align=center,font=\footnotesize] at (0,-2.8) {one write per\\pre-key rotation};
  \node[font=\small] at (5.4,1.2) {\textbf{Continuous mode}};
  \node[box] (k0) at (4.4,0.4)  {$k_0$};
  \node[box] (k1) at (4.4,-0.7) {$k_1$};
  \node[box] (k2) at (4.4,-1.8) {$k_2$};
  \draw[->] (k0) -- (k1); \draw[->] (k1) -- (k2);
  \node[box] (f) at (6.6,-0.7) {fingerprints};
  \draw[->,dashed] (k1) -- (f);
  \node[align=center,font=\footnotesize] at (5.4,-2.8) {one write per\\ratchet epoch};
\end{tikzpicture}

\caption[Session mode versus continuous mode.]{Session mode derives labels from the initial secret
$\sharedS$ and registers the identity key once per pre-key rotation. Continuous mode derives a
label and fingerprint from the pre-update state $k_t$ at every epoch and authenticates the public
ratchet key $T_t$. The verifier advances only after a \textsc{Consistent} result. A \textsc{Fork} or
\textsc{NotPosted} result leaves the state unchanged.}
\label{fig:design-modes}
\end{figure}

The two modes differ in authentication coverage and cost (\cref{fig:design-modes}). \emph{Session
mode} authenticates the identity key with one write per pre-key rotation. \emph{Continuous mode}
authenticates each ratchet epoch with one write per epoch and detects the post-compromise active
man-in-the-middle that is ACKA's principal security guarantee, at a substantially higher log-write
rate. The two are modes of one system, sharing the prefix tree, the log, the STH, and the warning
channel, and \cref{ch:eval} compares their costs. Two limitations bound the continuous-mode
prototype. First, in the prototype the chained secret is advanced deterministically from the shared
secret. A production binding would replace this simulated epoch input with Signal's live Double
Ratchet chain key $ck$, obtained together with the public ratchet key $T_t$ at a DH-ratchet-step
integration point, which is a distinct hook from the session-establishment capture of the initial
shared secret $\sharedS$ (\cref{ch:discussion}). Second, the construction inherits ACKA's synchronous
epoch model, in which the label chain advances in lockstep on both sides, so a ratchet message lost
before the receiver processes its epoch would leave the two chains permanently divergent and every
subsequent check failing. The Double Ratchet itself tolerates such loss through skipped-message keys.
A production binding would therefore need to carry explicit epoch indices and a bounded skipped-epoch
window in the same manner (\cref{ch:discussion}). Cross-implementation test vectors also lock the Rust
and TypeScript CEA constructions (\cref{sec:tests-xlang}).

% ============================================================================
\FloatBarrier
\section{Windowed Consistency Proofs}\label{sec:design-windowed}

The windowed consistency proof meets the catch-up requirement of \cref{ssec:bg-implications}, removing the linear cost
of offline catch-up, and so achieves G2 (Consistency Across Offline Windows).

\subsection{The Catch-Up Problem and the Bound}\label{ssec:design-catchup}

A client returning after $W$ offline epochs must verify that the current directory state
is an append-only extension of the last state it previously trusted. A
straightforward approach verifies one consistency proof for every missed epoch, giving $W$
separate proofs and an overall verification cost of $O(W \cdot \log n)$. Successive
consistency proofs contain substantial overlap, making repeated verification redundant. Instead, the
client requests a single consistency proof spanning its previously trusted log size $m$
and the current log size $n$, returned together with the latest signed tree head.
RFC~9162~\cite{rfc9162} guarantees that such a proof has size
\begin{equation}\label{eq:design-proofsize}
  |\pcons(m, n)| \le \lceil \log_2 m \rceil + \lceil \log_2 n \rceil \quad\text{hashes},
\end{equation}
growing only logarithmically with the current log size $n = m + W$ rather than linearly
with the number of intermediate epochs (\cref{fig:design-windowed}). The bound in
\cref{eq:design-proofsize} limits the window's cost. The client verifies the proof against the signed
root of the new tree head and accepts the new state only as a signed, append-only extension of its
trusted state. The warning-slot lookup is
performed during the same synchronization workflow and so does not change the asymptotic
complexity. \Cref{lem:windowed-size} states the windowed-proof size, and \cref{app:algorithms} gives the
pseudocode for all eight algorithms.

\begin{figure}[H]
\centering
% fig-design-windowed.tex — extracted from ch4-design.tex
% Label: \label{fig:design-windowed}
\begin{tikzpicture}[>=Stealth,font=\footnotesize]
  \draw[thick] (-0.3,0)--(8.3,0);
  \foreach \x/\lab in {0/$m$,1.6/{},3.2/{},4.8/{},6.4/{},8/$n$}{
    \filldraw (\x,0) circle (1.6pt) node[below=3pt]{\lab};}
  \node[left,align=right,gray] at (-0.5,0){signed\\tree heads};
  % naive: one consistency proof per missed epoch (upper arcs)
  \foreach \a/\b in {0/1.6,1.6/3.2,3.2/4.8,4.8/6.4,6.4/8}{
    \draw[->,red!70!black] (\a,0.12) to[bend left=55] (\b,0.12);}
  \node[red!70!black] at (4,1.55){na\"ive: $W$ proofs, cost $O(W\log n)$};
  % windowed: single proof m -> n (lower arc)
  \draw[->,green!45!black,very thick] (0,-0.18) to[bend right=32] (8,-0.18);
  \node[green!45!black] at (4,-1.5){windowed: one proof $\pcons(m,n)$, cost $O(\log n)$};
\end{tikzpicture}

\caption[Windowed offline catch-up.]{Offline catch-up over $W=n-m$ epochs. The upper arcs show
the synthetic evaluation baseline, which verifies one consistency proof per missed epoch at
$O(W\log n)$. The lower arc shows the RFC~9162 windowed approach, which requests one
consistency proof between the trusted head at $m$ and the current head at $n$. Its proof size is
logarithmic in the log size.}
\label{fig:design-windowed}
\end{figure}

The key-transparency literature identifies a soundness risk when clients skip intermediate epochs.
CONIKS's monitoring cost is linear in the offline period precisely because a malicious provider could
insert a false binding at some epoch inside the window and revert it before the client returns. A
client that verifies only the consistency of the two endpoint heads and the current value at its own
label would miss such a transient substitution, which is why CONIKS concludes that inspecting only the
latest signed root is unsound (\cref{ssec:bg-vkd-academic}). That objection applies verbatim to any
windowed check, including this one. A transient in-window substitution is profitable to the adversary
only if it is \emph{used}. A peer who initiates or continues the conversation during the window
performs pinned-key lookups at the dual-index slots (\cref{ssec:design-detection}) and, on a mismatch,
writes to the shared warning slot. DUET preserves detection only when that warning commits. The
write-once policy blocks replacement through the honest interface. Once committed, the warning either
remains available for the returning client to read or its removal causes the transition audit to
reject the epoch transition. The separately specified persistence check can retain evidence of that
removal.
If the server censors the write (the A2 threat of \cref{ch:threat}), the distant-head proof still
establishes append-only consistency but cannot exclude every transient substitution. Windowed catch-up
therefore relies on the warning channel, although the warning channel does not depend on windowing.
Without such a channel, as in CONIKS, the same windowing would weaken detection exactly as CONIKS's
authors argue.

\subsection{Worked Parameterization}\label{ssec:design-param}

Consider a directory with $N = 10^7$ users, about $c = 5$ active conversations per user, and signed
pre-keys rotating every $r = 30$ days.
With two party-specific lookup labels per conversation, the $N \cdot c \approx 50$~million figure
counts user--conversation lookup entries (about 25~million conversations, each with two lookup slots).
Each conversation also allocates one shared warning label, normally unwritten. Amortized over the
illustrative 30-day rotation period, the lookup entries are about 1.67~million directory updates per
day, or roughly $69{,}400$ updates folded into each hourly epoch commit (the log itself
gains one leaf per epoch regardless of update volume, \cref{ssec:design-epoch}). A
client offline for a full rotation period faces a worst-case window of $W = r \cdot 24 =
720$ epochs. Under the naïve scheme this is 720 consistency proofs. Under the windowed
scheme it is one, of size $O(\log n)$ in the log size $n$. \Cref{ch:eval} evaluates this
parameterization.

% ============================================================================
\FloatBarrier
\section{Instantiating ACKA}\label{sec:design-acka}

ACKA composes continuous key agreement (CKA), continuous entity authentication (CEA), and an abstract
secure-log interface $\Log$, as described in \cref{sec:bg-acka}. \Cref{tab:design-acka} maps these
components onto their instantiation in this work.

\begin{table}[htbp]
\centering
\caption{Mapping the ACKA components onto their instantiation in this work.}
\label{tab:design-acka}
\footnotesize
\begin{tabularx}{\textwidth}{@{}l>{\raggedright\arraybackslash}p{3.2cm}>{\raggedright\arraybackslash}X@{}}
\toprule
ACKA component & Role in the framework & Instantiation in this work \\
\midrule
$\CKA$ & continuous key agreement & Signal's Double Ratchet, unmodified at the design level, simulated in the prototype \\
$\CEA$ & per-epoch labels and fingerprints & Continuous mode retains ACKA's per-epoch CEA label-and-fingerprint construction \\
$\Log$ & abstract append-only secure log, required to be $\Logsec$-secure & DUET's verifiable key directory, comprising the prefix tree, Merkle log, signed tree heads, consistency and exclusion proofs, the write-once raw-label policy, and per-epoch transition validation \\
\bottomrule
\end{tabularx}
\end{table}

DUET's prefix tree, Merkle log, signed tree heads, consistency and exclusion proofs, write-once
raw-label policy, and per-epoch transition validation instantiate ACKA's abstract secure-log
component $\Log$. ACKA's $\CEA$ derives per-epoch labels and fingerprints from chained state, and DUET
retains ACKA's per-epoch CEA label-and-fingerprint construction in continuous mode. DUET separately
adds conversation-derived dual-index lookup labels and a shared warning slot in session mode, the
in-band peer-warning channel ACKA lacks. The windowed consistency policy of \cref{sec:design-windowed}
adds the deferred, batched catch-up that ACKA's per-epoch verification does not provide. The
cross-vantage comparison of signed heads adds split-view equivocation detection beyond ACKA's
$\Logsec$ requirement. The result is an implemented and measured $\Log$ component, with separate
extensions for peer warning, deferred catch-up, and split-view detection.

% ============================================================================
\section{Design Decisions and Alternatives}\label{sec:design-decisions}

The design rests on a number of choices, each made against a specific alternative and,
where possible, with a deployed precedent. \Cref{tab:design-decisions} summarizes them, and the
paragraphs below expand the three that most affect the system's behaviour.

\begin{table}[htbp]
\centering
\caption{Principal design decisions.}
\label{tab:design-decisions}
\setlength{\extrarowheight}{2pt}
\small
\begin{tabularx}{\textwidth}{@{}>{\raggedright\arraybackslash}p{3.7cm}>{\raggedright\arraybackslash}X>{\raggedright\arraybackslash}p{3.3cm}@{}}
\toprule
Decision & Justification & Reference / precedent \\
\midrule
Sparse binary Merkle tree & uniform node semantics, cheap first-class non-inclusion for the warning slot, small audit surface & CONIKS~\cite{melara2015coniks}, IETF KeyTrans~\cite{ietf-keytrans-protocol} \\
ECVRF-EDWARDS25519-SHA512-TAI for global labels & publicly verifiable directory privacy, parity with Signal's deployed keytrans VRF & RFC~9381~\cite{rfc9381}, Signal keytrans~\cite{signal-kt-server} \\
SHA-256 with RFC~6962 domain tags & second-preimage separation of leaves and nodes, RFC~9162 alignment & RFC~6962/9162~\cite{rfc6962,rfc9162} \\
HKDF-SHA256 for dual-index labels & standard, X3DH-consistent KDF, PRF security underpins pre-use unpredictability & RFC~5869~\cite{rfc5869}, Signal X3DH~\cite{x3dh} \\
Two-tree (prefix + log) combined stack & one verifiable chain from entry to signed head & IETF KeyTrans draft~\cite{ietf-keytrans-protocol} \\
Hourly epochs (tunable) & balances proof freshness against commit cost & CONIKS~\cite{melara2015coniks} \\
Client-tunable window $W$ & lets each client trade catch-up frequency against per-check cost & this work \\
Write-once dual-index and CEA labels, audited per epoch & preserves the entry immutability required by ACKA's $\Logsec$ notion, makes warning deletion detectable and provable & ACKA~\cite{dowling2023acka}, IETF KeyTrans auditing~\cite{ietf-keytrans-protocol} \\
Truncate VRF output $\beta$ to 32 bytes & 256-bit label space, matches the prefix-tree label width & Signal keytrans~\cite{signal-kt-server} \\
\bottomrule
\end{tabularx}
\end{table}

\paragraph{Directory Structure.} The directory is a sparse binary Merkle tree rather
than an $n$-ary trie or a cryptographic accumulator. A binary tree gives uniform node
semantics: every internal node hashes two children under the RFC~6962 node tag. This
keeps both the client verifier and the audit surface small, and it makes
non-inclusion a first-class, cheap proof: an absent label resolves to a default-hash path,
which is exactly what the warning tripwire (\cref{ssec:design-detection}) relies on to
prove a slot empty. An accumulator would shorten membership proofs but make the
non-membership proofs the warning channel needs substantially more expensive, and an
$n$-ary trie would reduce the number of levels but widen every node's copath ($k-1$ siblings per
level) and complicate the encoding, while with path compression the transmitted proof is already
governed by the populated prefix depth, so binary branching keeps the simplest node semantics at no
asymptotic cost in proof size. This is the structure used by CONIKS and the IETF Key Transparency design
(\cref{ssec:bg-keytrans}).

\paragraph{Global-Label Derivation.} Global slots are addressed by an ECVRF over the user identifier
rather than a bare hash, for the directory-privacy reason given in \cref{ssec:design-structures}. The
design selects ECVRF-EDWARDS25519-SHA512-TAI (RFC~9381) because it matches Signal's deployed keytrans
VRF, avoiding a second, incompatible VRF suite and preserving compatibility with Signal's
key-transparency service. Its measured cost in this prototype is reported in \cref{ch:eval}.

\paragraph{Combined Prefix-Plus-Log Stack.} Authentication uses two trees, a prefix
tree for the current directory and an append-only log of its committed roots, rather
than a single structure. One tree alone would authenticate the present state but not its
history, leaving append-only violations undetectable. The log supplies exactly the
consistency proofs that windowed catch-up (\cref{sec:design-windowed}) and the
cross-vantage auditor G4 (Fork Detection) depend on. The cost is one extra log-inclusion path per
lookup, up to $32\lceil\lg E\rceil$ bytes, quantified in \cref{ch:eval}.

\Cref{tab:threat-goals} maps the four security requirements of \cref{sec:threat-goals} to their
associated threats and the mechanisms developed in this chapter.

\begin{table}[H]
\centering
\caption{Security requirements, associated threats, and supporting mechanisms.}
\label{tab:threat-goals}
\setlength{\extrarowheight}{2pt}
\footnotesize
\begin{tabularx}{\textwidth}{@{}ll>{\raggedright\arraybackslash}Xl@{}}
\toprule
Requirement & Associated threat & Security mechanism & Auditor? \\
\midrule
G1 key-substitution detection & A1, A3, A5 & A1 and A5 use self-checks and the dual-index warning.\newline A3 uses divergent labels and the global VRF-bound identity entry.\newline A5 also uses identity re-fetch and predecessor cross-signature. & no \\
G2 offline-window consistency & general rewriting & windowed consistency proofs & no \\
G3 directory privacy & enumeration / recon & VRF-derived labels & no \\
G4 fork detection & A4 & cross-vantage auditor & yes ($\geq 1$ honest) \\
\bottomrule
\end{tabularx}
\end{table}
